% Borrador de la sección 1.1 — Contexto y motivación (sec:intro_contexto).
% Ficha: fichas/01-introduccion.md, sección 1.1. Extensión objetivo: 3--3,5 páginas.
% Al integrar: reemplaza el \pendiente de la sección homónima en Chapters/01-Introduccion.tex.
% Claves citadas: sung2026global, nci2026melanoma, loria2020cutaneous, vestergaard2008dermoscopy,
% kittler2002diagnostic, tschandl2019comparison, freeman2026global, sinsky2016allocation,
% esteva2017dermatologistlevel, magalhaes2024systematic, kreouzi2024learning, jones2022intelligence,
% daneshjou2022disparities, groh2021neural, riccilara2023dataset, radford2021learning, liu2023visual,
% li2023llavamed, zhou2023skingpt4, yan2025derm1m, hu2021lora, yilmaz2025resourceefficient.
% Figura: Figures/Capitulos/01-planteo-problema.tex (fig:intro_dos_tareas).

\section{Contexto y motivación}
\label{sec:intro_contexto}

El cáncer de piel es uno de los tipos de cáncer más frecuentes del mundo: las estimaciones de GLOBOCAN para
2024 le atribuyen 1\,138\,314 casos nuevos de cáncer de piel no melanoma, el quinto tumor más frecuente del
planeta, y otros 337\,969 casos de melanoma~\cite{sung2026global}. El melanoma es así menos de una cuarta
parte de los diagnósticos, pero las muertes se reparten casi por partes iguales entre los dos grupos, 58\,398
por melanoma y 58\,721 por el resto de los cánceres de piel, de modo que esa fracción menor de los casos
concentra prácticamente la mitad de la mortalidad. Esa asimetría entre frecuencia y letalidad es la que da
sentido clínico al problema: la mayoría de las lesiones que llegan a una consulta son benignas o no
representan un peligro real y el valor de examinarlas está en identificar la minoría que no lo es.

El peso de la detección temprana descansa en que el desenlace del melanoma depende del estadio en que se lo
diagnostica, de modo que adelantar el momento del diagnóstico no mejora el tratamiento sino el pronóstico. La
supervivencia relativa a cinco años —la proporción de pacientes que sobrevive una vez descontada la
mortalidad esperada en la población general— es de 100.0\,\% cuando el melanoma se detecta localizado, baja a
76.0\,\% cuando ya alcanzó los ganglios regionales y a 34.0\,\% cuando hay metástasis a
distancia~\cite{nci2026melanoma}. En esa misma serie, 77\,\% de los casos se diagnostica localizado, 10\,\%
con extensión regional y 5\,\% con metástasis a distancia: el margen de mejora está en ese último sexto de
los diagnósticos, y no en la mayoría que ya se detecta a tiempo.

En Argentina, el Registro Argentino de Melanoma Cutáneo reunió 10\,199 casos entre 2002 y 2018 y documentó
diferencias regionales en las características de la enfermedad~\cite{loria2020cutaneous}; es la fuente
disponible para caracterizar el problema en el país y la referencia con la que se contrasta cualquier
afirmación local.

La vía humana para ese diagnóstico no está resuelta. La dermatoscopía, que examina la lesión con aumento e
iluminación controlada, mejora la exactitud diagnóstica respecto de la inspección a ojo desnudo, y la mejora
se verifica en la práctica clínica y no solo en condiciones de estudio~\cite{vestergaard2008dermoscopy}. Pero
no es un beneficio que venga con el instrumento: depende de la formación de quien lo usa, y en examinadores
sin entrenamiento específico la dermatoscopía no aporta ventaja sobre la inspección
simple~\cite{kittler2002diagnostic}.

A la dependencia de la experiencia se suma la dispersión entre observadores. Un estudio de lectores de gran
escala enfrentó a 511 lectores humanos con 139 algoritmos sobre el mismo material dermatoscópico y encontró
un rango amplio de desempeño entre los lectores, asociado a su nivel de
formación~\cite{tschandl2019comparison}. Esa dispersión tiene dos consecuencias para este trabajo: fija la
referencia contra la cual se mide cualquier sistema automático, y advierte que una etiqueta producida por un
único lector no constituye una referencia estable.

El tercer factor es de escala y no de dificultad. Un relevamiento de la densidad de dermatólogos en 158
países estimó una media de 2.66 especialistas por cada 100\,000 habitantes, con un rango que va de 0.37 en
los países de ingreso bajo a 5.05 en los de ingreso alto~\cite{freeman2026global}. Donde el especialista es
escaso, la primera lectura de una lesión queda a cargo de quien no tiene formación dermatoscópica, que es
justamente el caso en el que el instrumento no aporta ventaja. A la escasez se agrega el costo de tiempo de
la documentación: la medición directa del uso del tiempo médico en la práctica ambulatoria atribuye
49.2\,\% de la jornada de consultorio a tareas de registro y trabajo de escritorio, más entre una y dos horas
fuera del horario de atención~\cite{sinsky2016allocation}. Redactar el informe de una lesión es parte de esa
carga, y es una tarea distinta de la de emitir el diagnóstico.

Sobre ese cuadro se instaló el aprendizaje profundo. Ya en 2017 se mostró que una \gls{cnn} entrenada sobre
un conjunto grande de imágenes de lesiones cutáneas alcanzaba un desempeño comparable al de dermatólogos
certificados~\cite{esteva2017dermatologistlevel}, y en 2019 el estudio de lectores ya citado encontró que los
mejores algoritmos superaban a los expertos humanos sobre el mismo tipo de
material~\cite{tschandl2019comparison}. Ambos resultados son anteriores a la adopción de los
\emph{transformers} en visión por computadora y se toman acá como punto de partida, no como el desempeño
alcanzable hoy. Las revisiones sistemáticas posteriores registran exactitudes altas en condiciones
controladas~\cite{magalhaes2024systematic,kreouzi2024learning}, mientras los trabajos que examinan el uso de
estos sistemas en atención primaria y comunitaria plantean una pregunta distinta: si ese desempeño se
conserva fuera de las condiciones de lectura especializada~\cite{jones2022intelligence}.

Ese estado actual deja dos huecos concretos. El primero es la forma de la salida. Un clasificador entrega una
categoría, o una distribución sobre categorías, y no un texto que describa la lesión y sostenga esa
categoría; la parte del trabajo clínico que consume tiempo queda sin asistencia. El segundo es la
generalización. La evaluación de sistemas de \gls{ia} dermatológica sobre conjuntos de imágenes curados para
cubrir el rango de fototipos muestra caídas de desempeño en los tonos de piel menos representados en los
datos de entrenamiento, que son los que dominan los conjuntos
públicos~\cite{daneshjou2022disparities,groh2021neural}. Por eso la validación externa de este trabajo se
hace sobre HIBA, un conjunto de imágenes de lesiones cutáneas recolectado en Argentina con el propósito
explícito de evaluar herramientas de \gls{ia} en esa población~\cite{riccilara2023dataset}.

Con ese encuadre, el problema que aborda el trabajo se puede precisar. El diagnóstico dermatológico asistido
involucra dos tareas y no una (Figura~\ref{fig:intro_dos_tareas}): asignar a la imagen una categoría
diagnóstica, y redactar un texto que describa la lesión y sostenga esa categoría. \emph{Detección} designa acá la tarea binaria de decidir si una
lesión es un melanoma, derivada de la clasificación multiclase agrupando sus categorías; no designa la
localización de la lesión dentro de la imagen ni su segmentación, que son problemas distintos y quedan fuera
del trabajo. Un reporte correcto exige entonces dos condiciones simultáneas: que el diagnóstico que afirma
sea el acertado, y que la descripción corresponda a la imagen que se está mirando. Las dos condiciones son
independientes entre sí, y esa independencia es lo que obliga a evaluarlas por separado: un texto fluido y
bien estructurado puede sostener un diagnóstico equivocado, y un diagnóstico acertado puede ir acompañado de
una descripción que no se corresponde con la lesión.

% Figura 1.1. Las dos tareas del diagnóstico dermatológico asistido.
% Se usa en sec:intro_contexto. El diagrama del pipeline completo (sec:met_pipeline) es la versión
% detallada de este mismo esquema: si uno de los dos cambia, hay que revisar el otro.

\begin{figure}[htbp]
    \centering
    \begin{tikzpicture}[x=1mm, y=1mm]

        \node[figentrada, text width=27mm] (img) at (12, 0) {Imagen\\dermatoscópica};

        \node[figtarea, text width=36mm] (t1) at (58, 14)
            {\textbf{Clasificar}\\asignar una categoría};
        \node[figtarea, text width=36mm] (t2) at (58, -14)
            {\textbf{Describir}\\redactar el texto\\de la lesión};

        \node[figsalida, text width=34mm] (cat) at (114, 14)
            {Categoría diagnóstica\\[2pt]{\footnotesize de ella se deriva la\\detección de melanoma}};
        \node[figsalida, text width=34mm] (txt) at (114, -14)
            {Reporte clínico\\[2pt]{\footnotesize el texto que\\describe la lesión}};

        \draw[figflecha] (img.east) -- ++(6, 0) |- (t1.west);
        \draw[figflecha] (img.east) -- ++(6, 0) |- (t2.west);
        \draw[figflecha] (t1.east) -- (cat.west);
        \draw[figflecha] (t2.east) -- (txt.west);

        \begin{scope}[on background layer]
            \node[figgrupo={lo que un modelo visión--lenguaje promete resolver con un solo modelo},
                  fit=(t1) (t2)] {};
        \end{scope}

    \end{tikzpicture}
    \caption[Las dos tareas del diagnóstico dermatológico asistido]{Las dos tareas del diagnóstico
    dermatológico asistido. De una misma imagen se espera una categoría diagnóstica y un texto que describa
    la lesión, y la detección de melanoma es la tarea binaria que se deriva de esa clasificación. Un reporte
    es correcto solo si se cumplen las dos condiciones a la vez: que la categoría sea la acertada y que la
    descripción corresponda a la imagen. Las dos condiciones son independientes y de ahí que se evalúen por
    separado.}
    \label{fig:intro_dos_tareas}
\end{figure}


La promesa de un \gls{vlm} es resolver las dos tareas con un único modelo, y hay una razón técnica concreta
para esperarlo. El punto de partida es el entrenamiento conjunto de un codificador de imágenes y un
codificador de texto sobre pares imagen--descripción, que ubica a ambas modalidades en un espacio de
representación común: una imagen se puede comparar con una frase, y la clasificación se vuelve la búsqueda
de la descripción textual que mejor le corresponde, sin un conjunto de etiquetas fijado de
antemano~\cite{radford2021learning}. Sobre esa base, los modelos generativos multimodales agregan un modelo
de lenguaje que produce texto libre condicionado a la imagen~\cite{liu2023visual}, y su adaptación al
dominio biomédico mostró que el procedimiento se traslada a imágenes médicas~\cite{li2023llavamed}.

En dermatología esos modelos ya existen. Un sistema diagnóstico interactivo acopló un modelo de lenguaje
visual a imágenes de lesiones cutáneas para sostener un intercambio en lenguaje
natural~\cite{zhou2023skingpt4}, y trabajos más recientes construyeron conjuntos de pares imagen--texto de
escala millonaria alineados con la ontología clínica del área, sobre los que se entrenaron modelos de
representación conjunta específicos del dominio~\cite{yan2025derm1m}. Que existan es lo que vuelve razonable
la pregunta del trabajo: si un único VLM puede a la vez clasificar la lesión y redactar su reporte, o si las
dos tareas siguen exigiendo componentes distintos y encadenados.

Queda un eje que atraviesa todo lo que sigue: el hardware. Lo que se reporta acá corre en una única GPU de
consumo de 12--24\,GB de memoria, la clase de equipamiento que hay en un laboratorio universitario o en un
servicio clínico, y no en un centro de cómputo con decenas de aceleradores. La restricción no es una
limitación a disculpar sino una condición de diseño, y decide qué modelos entran en la comparación, con qué
resolución y tamaño de lote se los entrena y cuántas configuraciones se alcanzan a explorar. Lo que vuelve
manejable el ajuste fino (\emph{fine-tuning}) en ese presupuesto es \gls{lora}, que mantiene congelados los
pesos del modelo pre-entrenado y entrena en su lugar un conjunto reducido de matrices de bajo rango, con lo
que la memoria requerida deja de escalar con el tamaño total del modelo~\cite{hu2021lora}; hay trabajos que
obtienen por esa vía un modelo visión--lenguaje dermatológico con recursos de ese
orden~\cite{yilmaz2025resourceefficient}. Sobre este encuadre —dos tareas, un tipo de modelo que promete
resolver ambas y un presupuesto de cómputo acotado— se formulan las preguntas que el trabajo responde.
